\documentclass[10pt,a4paper]{amsart}
\usepackage[margin=19mm]{geometry}
\usepackage{amsmath,amssymb,mathtools}
\usepackage[scr=rsfs,cal=euler]{mathalfa}
\usepackage{xfrac}
\usepackage[unicode,hidelinks,bookmarks=false]{hyperref}
\newcommand{\ie}{i.e.\ }
\newcommand{\cf}{cf.\ }
\newcommand{\resp}{respectively\ }
\newcommand{\Subsection}{\subsection}
\providecommand{\nobreakdash}{\nobreak\hbox{-}}
\setlength{\emergencystretch}{2em}
\multlinegap0pt
\usepackage{float}
\usepackage[shortlabels]{enumitem}
\usepackage{amssymb}
\usepackage{tcolorbox}
\usepackage{cancel}
\usepackage[normalem]{ulem}
\usepackage{tikz}
\newtheorem{lemma}{Lemma}
\newtheorem{theorem}{Theorem}
\newcommand{\M}{\mathcal{M}}
\title{Braiding Vineyards}
\author{Erin Chambers, Christopher Fillmore, Elizabeth Stephenson, Mathijs Wintraecken}
\date{Source: arXiv:2504.11203v2 (CC BY 4.0)}
\begin{document}
\maketitle

\noindent\emph{Source and licence.} Braiding Vineyards. Version arXiv:2504.11203v2 (CC BY 4.0), distributed by arXiv under the Creative Commons Attribution 4.0 International licence, http://creativecommons.org/licenses/by/4.0/, as recorded in the arXiv licence metadata for this exact version and shown on the abstract page https://arxiv.org/abs/2504.11203v2. Full licence notice: \texttt{licenses/arxiv-2504.11203-CC-BY-4.0.txt}.\par\smallskip

\noindent\textit{Editorial scope.} Counted unit: the article's theorem th:Morse (source0.tex lines 1293--1295). The article's complete original proof of it is delivered with it (source0.tex lines 1297--1458, one \texttt{proof} environment of 162 lines). No mathematics is written, repaired or re-derived: the statement and the proof are the article's own lines, in the article's own notation, and no retained line is substituted or re-flowed.  Retained uncounted as the source-local context the proof needs: lines 685--693.  One declared adaptation: exactly 2 ancillary strings inside the retained spans are omitted (image-inclusion commands and, where the source repeats a label, the repeated label definitions) and no other line differs from the source; the figure environments, with their placement, captions and labels, are kept, and the package records the omitted strings in fidelity receipts bound to the affected bodies.  No retained line contains a citation, so the wrapper carries no bibliography and no bibliography entry.  The retained lines contain the article's own diagram code, which the wrapper typesets with the standard tikz packages; no image file is delivered and the package declares no asset dependency.  Editorial formatting only, in addition: an AMS \texttt{amsart} preamble that restates the article's own declarations and macros used by the retained lines, namely the environments \texttt{lemma}, \texttt{theorem} on separate counters, together with the standard packages those lines need.  The retained environments print in this document as Lemma 1, Theorem 1 in order of appearance, whereas the article prints the same items with the numbering of its own full document.  The complete native source of the article as distributed in the arXiv e-print for this exact version is bundled unchanged as \texttt{sources/176577/source0.tex}.
\begin{lemma}[clustering of critical points] \label{Lem:Step2} Let $B$ be an $(\epsilon,R)$-embedded closed braid and suppose that $\M$ is its $(l+1)$-dimensional $\alpha$-offset, with braid angle $\theta_B$. If $p \in \mathbb{R}^d$ satisfies $d (p, C_h(0,R)) \leq \eta$, $C_h(0,R)$ is parametrized by $s(\theta)$, and the closest point $p'$ of $p$ on $C_h(0,R)$ satisfies $p' =s(0)$, then the function $x \mapsto d_\mathbb{E}(x, p)|_\M$ has no critical points at the closest point $\beta (\theta)$ on $B$ to $x$ to as long as  
\begin{align}  \frac{\theta}{2} + \theta_B + \arcsin \frac{\epsilon} { 2 R \sin (\frac{\theta}{2})  -\eta } +\arcsin \frac{\eta} { 2 R \sin (\frac{\theta}{2})  }  <\frac{\pi}{ 2} .
\label{eq:BoundL1}
\end{align} 
This implies in particular that there is no topological change as long as \eqref{eq:BoundL1} is satisfied. 
%
%
%
\end{lemma} 

\begin{theorem}[Morse]\label{th:Morse} 
Suppose $f$ is a smooth real-valued function on $\M$, $a < b$, $f^{-1} [ a , b ]$ is compact, and there are no critical values between $a$ and $b$. Then $\M^a$ is diffeomorphic to $\M^{b}$, and $\M^{b}$ deformation retracts onto $\M^{a}$. 
\end{theorem}

\begin{proof}[of Lemma \ref{Lem:Step2}] 
We write $\beta (\theta)$ for a parametrization of $B$ according to the angle of the circle $C_h(0,R)$. We parametrize the circle as $s(\theta) =( R \sin \theta,R\cos \theta, 0 , \dots ,0 )$. Using this notation, we note that the gradient of $d(y, p)_\M$ is zero at $y=x$ if for its closest point on $B$, that is, $\beta(\theta)$, we have that 
$ \langle {\beta}',\beta- p  \rangle =0 $, where $\beta'$ denotes the derivative of $\beta$ with respect to $\theta$. This is equivalent to 
\[ \angle  {\beta}',\beta- p  = \pi/2.
\]
By assumption we have that 
\[  \angle  {\beta}' , s' \leq \theta_B.
\]
Moreover, because $| \beta(\theta) - s(\theta) | \leq \epsilon $, we also have 
\[ \sin (\angle \beta -p , s- p)   \leq  \frac{\epsilon} { |s -p| }, 
\]
as can be seen from Figure \ref{fig:Obvious}. 
%

\begin{figure}[h!]
    \centering
    
    \caption{The angle estimate for $\angle \beta -p , s- p $. }
    \label{fig:Obvious}
\end{figure}


%
%
%
%
%
%
%
%
%
%
%
%
%
%
%
%
%
%
%
%
%
%
%
%
%
%
%
%
%
%
%
%
%
%
%
%
%
%
%
%
%
%
%
%
%
%
%
%
%
%
%
%
%
%
%
%
%

Using the triangle inequality of angles (or points on the sphere) we find that 
\begin{align} 
|\angle  ({\beta}',\beta- p)  -  \angle  ({s}',s- p) | \leq  \theta_B + \arcsin \frac{\epsilon} { |s -p| } . 
\label{eg:TriangleIneq1}
\end{align} 
%

Let us now write $d= | s- p|$, $p'$ for the closest point projection of $p$ on $C_h(0,R)$
%
and $\omega = \angle  ({s}',s- p')$. 
By reparameterization we can assume that $p'= s(0)$. With this assumption we see by the construction in Figure \ref{fig:AngleEst3} that  $\omega=\angle  ({s}',s- p') = \frac{\theta}{2}$. Moreover we have that $|p'- s | = 2 R \sin (\frac{\theta}{2}) $
\begin{figure}[h!]
    \centering
    
    \caption{The construction for the angle $\omega= \angle  ({s}',s- p')$. }
    \label{fig:AngleEst3}
\end{figure}
By the same argument as that given in Figure \ref{fig:Obvious} we have  
\[  \sin \angle  (s-p',s- p)  \leq \frac{\eta}{|s-p'|},
\]
so that together with \eqref{eg:TriangleIneq1} we find that
\begin{align} 
\left |\angle  ({\beta}',\beta- p)  - \frac{\theta}{2}  \right|& \leq  \theta_B + \arcsin \frac{\epsilon} { |s -p| } +\arcsin \frac{\eta} { |s -p'| } 
\nonumber
\\
&\leq \theta_B + \arcsin \frac{\epsilon} { |s -p'| -\eta } +\arcsin \frac{\eta} { |s -p'| } 
\tag{by the triangle inequality} 
\\
&\leq \theta_B + \arcsin \frac{\epsilon} { 2 R \sin (\frac{\theta}{2})  -\eta } +\arcsin \frac{\eta} { 2 R \sin (\frac{\theta}{2})  } 
\nonumber
\end{align} 
This means that if 
\[ \frac{\theta}{2} + \theta_B + \arcsin \frac{\epsilon} { 2 R \sin (\frac{\theta}{2})  -\eta } +\arcsin \frac{\eta} { 2 R \sin (\frac{\theta}{2})  }  <\frac{\pi}{ 2} 
\] 
then  $\angle  ({\beta}',\beta- p) \neq \pi/2 $ and hence there are no critical points. The fact that there is no topological change follows from Theorem \ref{th:Morse}. 
%
%
%
%
%
%
%
%
%
%
%
%
%
%
%
%
%
%
%
%
%
%
%
%
%
%
%
%
%
%
%
%
%
%
%
%
%
%
%
%
%
%
%
%
%
%
\end{proof} 

\end{document}
